% -*- coding: utf-8 -*-
\chapter{Problema, contexto y utilidad real}

Este Trabajo de Fin de Máster diseña, implementa y evalúa un asistente documental basado en Generación Aumentada por Recuperación (RAG) sobre documentación agrícola y normativa. El problema principal no es ``aplicar inteligencia artificial al sector agrario'' en abstracto, sino \textbf{reducir alucinaciones factuales}: respuestas fluidas que inventan importes, plazos, códigos de ayuda o citas no presentes en el expediente (Huang et al., 2025; Ji et al., 2023). La pregunta operativa es si anclar la generación a fragmentos recuperados \textbf{mitiga} ese riesgo frente al mismo LLM sin recuperación, en un corpus PAC/POSEI y de sanidad vegetal (Lewis et al., 2020).

\section{Justificación del proyecto}

Los LLM han hecho trivial consultar texto en lenguaje natural. En dominios regulados esa facilidad es también un riesgo: una cifra plausible pero falsa puede orientar mal a un asesor o a un gestor de ayudas. La literatura trata la alucinación como barrera de fiabilidad, no como un defecto estético. En este trabajo se distingue:

\begin{itemize}
  \item \textbf{Alucinación intrínseca:} la salida contradice el contexto que el modelo tiene delante.
  \item \textbf{Alucinación extrínseca:} la salida afirma hechos que no están en el contexto (o no hay contexto): artículos, excepciones o códigos verosímiles pero nulos.
\end{itemize}

En normativa agraria (PAC, POSEI, GIP, cuarentena) predomina el segundo tipo: el modelo rellena huecos con patrones de boletín. Un LLM aislado no consulta el expediente; muestrea tokens a partir de pesos congelados en el preentrenamiento. Actualizar ese conocimiento paramétrico por reentrenamiento es costoso e inflexible cuando cambian convocatorias o textos consolidados.

RAG desacopla razonamiento y conocimiento: recupera fragmentos de un corpus externo y condiciona la generación a esa evidencia (Lewis et al., 2020; Gao et al., 2024). No es un mecanismo que ``lleve las alucinaciones a cero''; es una hipótesis a medir. Este TFM aporta un caso aplicado ---corpus real de asesoría canaria, banco de preguntas con fuente, comparación pareada con/sin RAG--- frente a benchmarks genéricos que no contienen POSEI ni registros fitosanitarios.

La utilidad del sistema no se define por la fluidez, sino por la \textbf{trazabilidad}: cada respuesta debe poder señalar documento, sección (y página cuando el parser la conserva) y fragmento recuperado. Sin esa cadena, el prototipo no sería distinguible de un chat genérico.

\section{Contexto}

El usuario tipo es quien consulta boletines, medidas de ayuda, condicionalidad y fichas de cultivo: documentos largos, con tablas y con vigencia temporal. El prototipo AgroPS opera en local (Docker Compose, Ollama) para que el corpus no salga a un servicio comercial; esa decisión es de ingeniería del MVP, no el objeto experimental. No se evalúan data lakes, cadenas de suministro, deepfakes ni productos tipo Watson: quedan fuera del alcance porque no intervienen en la comparación LLM autónomo frente a RAG.

\section{Presentación del problema y de la solución}

Un LLM sin recuperación puede responder con corrección lingüística y error factual. AgroPS implementa un \textbf{RAG textual modular}: ingesta a texto (PDF/OCR cuando procede), índice denso (\texttt{pgvector}) más índice léxico (BM25), fusión RRF y generación local. El chat muestra fuentes. El modo agéntico (LangGraph) existe en el producto; la evaluación cuantitativa principal mide el recuperador híbrido con top-$k$ fijo. GraphRAG y visión end-to-end no se evalúan.

El trabajo sigue un ciclo de ingeniería aplicada: preparar el corpus, fragmentarlo, indexarlo, generar con citación y comparar, sobre un banco versionado, las respuestas con y sin recuperación.

\section{Utilidad real}

En lugar de buscar a mano entre PDF, el usuario formula una pregunta y obtiene una respuesta con fragmentos. Eso reduce el tiempo de localización y, sobre todo, permite \textbf{auditar} de dónde sale cada dato. En un entorno donde la normativa cambia por campaña, actualizar el índice es más realista que reentrenar el modelo. La citación no garantiza cumplimiento jurídico; sí hace visible el fallo cuando el retriever trae el documento vecino y el generador inventa un código (Capítulo~\ref{ch:evaluacion}).

\section{Estructura del trabajo}

La memoria se organiza del siguiente modo:
\begin{enumerate}
  \item Introducción: problema de alucinación en documentación agraria.
  \item Objetivos acotados y planificación del MVP.
  \item Estado del arte crítico (RAG, alucinaciones, métricas, dominios regulados).
  \item Metodología: corpus, banco de preguntas, variables y protocolo.
  \item Arquitectura e implementación del RAG textual modular.
  \item Despliegue del prototipo contenedorizado.
  \item Evaluación: LLM sin RAG frente a LLM con RAG, y discusión.
  \item Conclusiones, limitaciones y trabajo futuro.
\end{enumerate}
